\section{Agile Projects and New Work}


\begin{frame}{`Stacey'-Matrix}
       \showmytikz{../shared/tikz/tikz-PM-stacey-matrix-requirements}
        \end{frame}

\begin{frame}{`Stacey'-Matrix}        
       \showmytikz{../shared/tikz/tikz-PM-stacey-matrix-stakeholders}
        
\end{frame}

\begin{frame}{VUCA}
    \begin{tabular}{ll}
    \begin{minipage}{4cm}
        \begin{description}
            \item[V]olatility
            \item[U]ncertainty
            \item[C]omplexity
            \item[A]mbiguity
        \end{description}
    \end{minipage}
         &  
    \begin{minipage}{6cm}
        \begin{description}
            \item[V]ision
            \item[U]nderstanding
            \item[C]larity
            \item[A]gility
        \end{description}
    \end{minipage}
\end{tabular}
    \mytext{
\medskip

Sources e.g.~\cite{AkkayaBuelent2022AMaV}
    }
\end{frame}


\begin{frame}{Culture vs. Strategy: Unternehmens`philosophie'}
    
    \mytext{
    \begin{quote}
        Plans are worthless, but planning is everything.
        
        ---Attributed to Dwight D.\ Eisenhower
    \end{quote}
    
    \medskip
    
    \begin{quote}
        `Culture eats strategy for breakfast'
        
        ---Attributed to Peter \noun{Decker}
    \end{quote}
    
    \medskip
    
    culture $\neq$ civilisation
    }
\end{frame}



\begin{frame}{Waterfall Management}
    
    \mytext{
    \begin{itemize}
        \item hierarchic
        \item pre-defined goals
        \item goals must be known!
        \item means must be known!
        \item tight control of subordinates
        \item inflexible
    \end{itemize}    
    Roles:
    \begin{itemize}
        \item project sponsor
        \item project manager
        \item subordinates
    \end{itemize}
    }
\end{frame}




\begin{frame}{Agility}
    
    \mytext{
    Paradigm shift: \fbox{waterfall, hierachy} $\rightarrow$ \fbox{agile, scrum}

    \medskip
    
    Change of business culture and organisation
    
    Project manager $\rightarrow$ product owner, scrum master
    
    Product backlog, sprints, kanban, daily huddle, scrum...
    
    Design Thinking
    
    ...
    
    }
\end{frame}

\begin{frame}{Agility}
    
    \mytext{
    
    \textbf{Agile Manifesto}

    \medskip

    \centering
    
    We are uncovering better ways of developing
    software by doing it and helping others do it.\\
    Through this work we have come to value:
    
    \medskip
    
    \textbf{Individuals and interactions} \small{over processes and tools}
    
    \textbf{Working software} \small{over comprehensive documentation}
    
    \textbf{Customer collaboration} \small{over contract negotiation}
    
    \textbf{Responding to change} \small{over following a plan}
    
    \medskip
    
    That is, while there is value in the items on
    the right,\\ we value the items on the left more.
    
    \vfill
    
    \href{https://agilemanifesto.org/}{https://agilemanifesto.org/}
        
    }
\end{frame}



\begin{frame}{Design Thinking Process - "Rapid prototyping"}

    \showmytikz{../shared/tikz/tikz-PM-design-thinking-wheel}

    % \begin{tikzpicture}[node distance=0.5cm,every node/.style={scale=0.8}]
    %     \usetikzlibrary{decorations.pathmorphing}
    %     \node[draw] (e) at (0,0) {empathise};
    %     \node[draw,right=of e] (d)  {\parbox{2cm}{define}};
    %     \node[draw,right=of d] (i)  {\parbox{1.1cm}{ideate}};
    %     \node[draw,right=of i] (p)  {\parbox{1.8cm}{prototype}};
    %     \node[draw,right=of p] (t)  {test};
    %     \draw[->] (e) -- (d);
    %     \draw[->] (d) -- (i);
    %     \draw[->] (i) -- (p);
    %     \draw[->] (p) -- (t);
    %     \draw[->,dashed] (d) to [bend left=45] (e); 
    %     \draw[->,dashed] (i) to [bend left=45] (d); 
    %     \draw[->,dashed] (p) to [bend left=45] (i); 
    %     \draw[->,dashed] (t) to [bend left=45] (p);
    %     \draw[->,dashed] (t) to [bend right=25] (e); 
    % \end{tikzpicture}
    
    
    \end{frame}

\begin{frame}{Design Thinking Process - "Rapid prototyping"}
    
    \mytext{
    
    d.school institute / Standford University
    
    \begin{description}
         \item[empathise] interviews, shadowing, seek to understand, non-judgemental
         \item[define] personas, role objectives, decisions, challenges, pain points
         \item[ideate] share ideas, all ideas worthy, diverge/converge, `yes and'-thinking, prioritise
         \item[prototype] mockups, storyboards, keep it simple, fail fast, iterate quickly
         \item[test] understand impediments, `what works?', role play, iterate quickly
    \end{description}
    }
\end{frame}


\begin{frame}{PDCA-Cylce}
    \showmytikz[caption={PDCA Cycle (Agile)},label=fig:pdca-cycle-agile]{../shared/tikz/tikz-PM-pdca-wheel-1}
    \mycomment{
     Plan - do - check - act/adjust
     
    }
\end{frame}



\begin{frame}{Agile vs.\ Waterfall / Traditional}
    
    \mytext{
    \begin{itemize}
        \item agile is \textbf{feature based,} not \textbf{requirement based}
        \item features are often fulfilled in parallel, while
        \item milestones / requirements are completed one after another 
        \item agile focuses on collaboration
        \item rather than on takeovers
        \item different information architecture
        \item \textit{closer }collaboration
    \end{itemize}
    }
\end{frame}

\begin{frame}{Agile vs.\ Waterfall / Traditional}
    
    \mytext{
    \begin{minipage}{0.7\textwidth}
        \begin{itemize}
            \item in agile, \textbf{time and budget often fixed}
            \item features / requirements not exactly clear
            \item quality assurance often more important than number of features
            \item results often measurable by usability, rather than technically
            \item features described by user stories
            \item humans more important than specs
        \end{itemize}
    \end{minipage}
    \vfill
    \cite{rubin2014essential}
    }
\end{frame}


\begin{frame}{Requirements and User Stories}
    
    \mytext{
    \begin{itemize}
        \item if you cannot (yet) describe the exact solution:
        \item describe the problem
        \item describe the features
    \end{itemize}

    To find requirements, write \textbf{user stories.}
    
    }
\end{frame}

\begin{frame}{Requirements and User Stories}
    
    \mytext{
    
    Not so good: 
    \par\medskip
    \begin{quote}
        Students want affordable shared office space.
    \end{quote}
    
    }
\end{frame}

\begin{frame}{Requirements and User Stories}
    
    \mytext{
    
    Better:\par\medskip
    \begin{quote}
        As a [\dots], I want [\dots] to achieve [\dots]
        \par\medskip    
        As a student,\\ I want a shared office space\\ to better prepare for my exams.
    \end{quote}
    
    }
\end{frame}

\begin{frame}{Requirements and User Stories}
    
    \mytext{
    
    Even better: concrete stories!
    \par\medskip
    Writing user stories helps you understand
    \begin{itemize}
        \item requirements
        \item must-haves and nice-to haves
        \item \textbf{user journeys!}
    \end{itemize}
    
    }
\end{frame}

\begin{frame}{Requirements and User Stories}
    
    \mytext{
    

    Alice is a $\dots$. In the situation $\dots$, she wants $\dots$. First, she hears about the product via $\dots$. Then, she logs into the website. Then $\dots$
    
    These are realistic stories, like little novellas.
    
    }
\end{frame}

\begin{frame}{Requirements and User Stories}
    
    \mytext{
    \textbf{Dieter} was born in Heidelberg. He is now a 22 years old student of Civil Engineering in Berlin. He wants to go on vacation. He has little experience or habit regarding his holidays; but he is seeking adventure, new people, and experience culture `from within'.
    
    }
\end{frame}

\begin{frame}{Requirements and User Stories}
    
    \mytext{
    
   \textbf{Ana} was born in Florence. She is now a 23 year old student of Comparative Linguistics in Rome. She shares a flat with two other students. In the summer, she regularly returns to Florence to her parents.
    
    }
\end{frame}

\begin{frame}{Requirements and User Stories}
    
    \mytext{
    
    \textbf{Jeff} is a 28 year old US-American entrepreneur with some experience with startups. On a back packing trip to Europe, he has met both Dieter and Ana. They wondered if there could be a way to make travelling for people their age easier, more affordable, and more fun than by booking hotel rooms. 
    
    }
\end{frame}

\begin{frame}{Requirements and User Stories}
    
    \mytext{
    
    \textbf{Carl} is a 64 year old clerk in a bank. (\dots)
    }
\end{frame}


\begin{frame}{Tasks, Stories, Epics}
        \showmytikz{../shared/tikz/tikz-tasks-stories-epics.tex}
    \end{frame}

\begin{frame}{Tasks, Stories, Epics}    
    \showmytikz{../shared/tikz/forest-themes-tasks-stories-epics}
        
\end{frame}

\begin{frame}{Planning Poker}
    \showmytikz{../shared/tikz/tikz-planningpokercards}
    \mycomment{
    "Planning poker" is a tool to identify the complexity\\ (cost, duration, man hours...) of tasks.
    
    Every team member gets a set of cards.
    
    
    
    Team members should be diverse
    
    \begin{itemize}
        \item programmers
        \item testers
        \item designers
        \item users
    \end{itemize}
    
    }
\end{frame}

\begin{frame}{Planning Poker}
    
    \mytext{
    
    \begin{itemize}
        \item Discuss work breakdown, ie. stories / tasks
        \item How "big" is this task?
        \item Card points will correspond to costs / time
        \item Put your card face down
        \item Outliers explain their reasoning
        \item Repeat
        \item Use to plan your progress
    \end{itemize}
    
    
    }
\end{frame}

\begin{frame}{Describe Process and Tools Rather Than Milestones}
    
    \mytext{
If you cannot plan milestones precisely,

    \begin{itemize}
        \item plan your method of working
        \item PDCA-cycles, agile methods
        \item bug tracking, ITC tools
        \item Quality assurance (e.g. pair programming)
    \end{itemize}
    }
\end{frame}




\begin{frame}{Scrum Roles}
    
    \mytext{
    \begin{minipage}{0.7\textwidth}
        \begin{description}
            \item[product owner] traditionally project sponsor; a management role; the customer; the customer within the company
            \item[scrum master] the moderator; in between project sponsor and the team; shields team from project owner direct intervention; a communicator; leads scrums, daily huddles etc.
            \item[team members] individual workers / coders / researchers; typically working in pairs \textbf{(pair programming)}
        \end{description}
    \end{minipage}
    }
\end{frame}



\begin{frame}{Product Backlog}
    
    \mytext{
    The PB is fundamentally a list of (required) features
    \begin{itemize}
        \item define required features
        \item break down into work packages
        \item prioritise
        \item list is called \textbf{product backlog}
    \end{itemize}
    
    While requirements may change during development,
    
    they should not change during sprint
    
    Work packages move to Kanban board
    
    }
\end{frame}

\begin{frame}{Product Backlog}
    \begin{tabular}{r l l}
        \toprule
        & Feature / Slide & assigned to \\
        \midrule
        1 & Idea presentation & Alice \& Ellen \\
        2 & Stakeholders & Bob \& Charlie \\
        3 & Governance / organisation & Bob \& Charlie \\
        4 & Users and user requirements & Doreen \\
        5 & Solution design concept & Alice \& Ellen \\
        6 & Work breakdown structure & Doreen \& Frank \\
        7 & Time table / schedule / Gantt diagram & Frank \& Gilles \\
        8 & Risk register / escalation management &  Gilles \& Herbert \\
        9 & Core resource requirements, sum up & Herbert \\
        \bottomrule
    \end{tabular}
    \mytext{
    
    Example:  \textbf{Project Presentation \& Governing Document}
    \par\medskip
    
    }
\end{frame}



\begin{frame}{Kanban Board}
    \showmytikz{../shared/tikz/tikz-PM-kanban-board-1}\end{frame}

\begin{frame}{Kanban Board}        \showmytikz{../shared/tikz/tikz-PM-kanban-board-2}\end{frame}

\begin{frame}{Kanban Board}        \showmytikz{../shared/tikz/tikz-PM-kanban-board-3}
\end{frame}

\begin{frame}{Sprints and Scrum}
    
    \mytext{
    \begin{itemize}
        \item rather than have milestones,
        \item have regular `meetings' (PDCA-cycles)
        \item a work period is called `sprint'
        \item work efficiently
        \item rather than to (artificial) deadlines
    \end{itemize}
    }
\end{frame}

\begin{frame}{Sprints and Scrum}
    
    \mytext{
    \begin{itemize}
        \item determine product backlog
        \item determine which features to attack first
        \item choose work packages for a sprint
        \item during sprint, have daily meetings aka scrums
        \item transparency (cloud based etc.)
        \item "definition of done"
    \end{itemize}
    }
\end{frame}


\begin{frame}{Sprint Planning}
    
    \mytext{
    \begin{minipage}{0.7\textwidth}
        \textbf{A sprint} is a \textbf{time-boxed effort}
        \begin{itemize}
            \item determine length of sprints 
            \item Select items from the product backlog to be achieved next
            \item assign to teams or individuals
            \item this is the \textbf{sprint backlog}
            \item e.g. two weeks (or one), one month
            \item review progress when each sprint is over
        \end{itemize}
    \end{minipage}
    }
\end{frame}

\begin{frame}{Sprint Planning}
    \showmytikz{../shared/tikz/tikz-PM-daily-scrum-cycle}
    \mycomment{
    
    During each sprint, have daily scrum
    \par\medskip
    
    
    }
\end{frame}

\begin{frame}{Sprint Planning}
    
    \mytext{
    
    In a daily scrum
    \begin{itemize}
        \item `briefly stick your heads together' \textbf{for 15 min only!}
        \item brief members on progress
        \item find help to remove impediments
        \item stand up meetings
        \item information other teams must have
        \item no discussion!
    \end{itemize}
    }
\end{frame}

\begin{frame}{Sprint Planning}
    \showmytikz{../shared/tikz/tikz-burndownchart}
    \mycomment{
    
    When sprint is over, have 
    \begin{itemize}
        \item Sprint review: present results
        \item Sprint retrospective: review processes
        \item Refine backlog (to include technical dept)
    \end{itemize}

    Repeat: next sprint planning
    \par\medskip
    Progress can be displayd on a \textbf{burndown chart.}
    
    
    }
\end{frame}

\begin{frame}{Cyclic Management}
    
    \mytext{
    
    \begin{itemize}
        \item circular thinking \textit{en vogue}
        \item nothing is without consequences
        \item learning is repetition
        \item refinement through repetition
        \item repetitions not identical
        \item iterative, increments, evolutionary
    \end{itemize}
    
    }
\end{frame}

\begin{frame}{Cyclic Management}
    % \resizebox{!}{0.8\textheight}{
    %     \begin{tikzpicture}
    %     \draw[->] (180:5) -- (0:5);
    %     \draw[->] (270:5) -- (90:5);
    %     \foreach [count=\i] \m in {Analysis, Evaluation, Development, Planning}
    %     \node at (-\i*90-135:5) {\i. \m};            
    %         \draw[->] [domain=0:25.5,variable=\t,smooth,samples=75]
    %     plot ({-\t r}: {0.007*\t*\t});
    % \end{tikzpicture}
    % }
    
    % \framebreak
    
    \showmytikz{../shared/tikz/tikz-PM-cyclic-management-wheel}
    \end{frame}

\begin{frame}{Cyclic Management}
    
    \mytext{
    
    \begin{itemize}
        \item many types of phases
        \item phases of different lengths
        \item cycles and sub-cycles
        \item increments of different sizes
        \item same resources / people per phase?
    \end{itemize}
    
    }
\end{frame}


\begin{frame}{Google's `Project Aristotle'}
    
    \mytext{
    ...looked at different (Google) teams
    
    trying to establish reason(s) for efficiency
    
    Findings:

    \begin{itemize}
        \item Teams work differently
        \item People work differently
    \end{itemize}
    
    Most important factor:\par
    \textbf{Psychological safety!}
    
    }
\end{frame}

\begin{frame}{Google's `Project Aristotle'}
    
    \mytext{
    \begin{minipage}{0.65\textwidth}
        \begin{description}
            \item[Psychological safety] team members feel safe to take risks and be vulnerable in front of each other
            \item[Dependability] team members get things done on time and meet (Google's) hight bar for excellence
            \item[Structure and clarity] team members have clear roles, plans, and goals
            \item[Meaning]  work is personally important to team members
            \item[Impact] team members think their work matters and creates change
        \end{description}
    \end{minipage}
    
     Source: \href{https://rework.withgoogle.com/print/guides/5721312655835136/}{https://rework.withgoogle.com/print/guides/5721312655835136/}    
        
    }
\end{frame}

\begin{frame}{Google's `Project Aristotle'}
    
    \mytext{
    \begin{itemize}
        \item Teams can be rule bound or chaotic
        \begin{itemize}
            \item Rules must fit characters and work types
            \item Efficiency correlates to psychological safety
        \end{itemize}
        \item Avoid a `no bad news' - environment
        \item Trust must be mutual and substantiated
        \item Trust is also based on commitment
    \end{itemize}        
        
    }
\end{frame}

\begin{frame}{Google's `Project Aristotle'}        
    \showmytikz{../shared/tikz/tikz-PM-project-aristotle-pyramid}
    
\end{frame}


\begin{frame}{Maslow's Hierarchy of Needs}
    \showmytikz{../shared/tikz/tikz-PM-maslow-hierarchy}
    \mycomment{
    
    \vfill
    Source: Maslow (1943) / commons
    }
\end{frame}


\begin{frame}{Tuckman's Phases of Team Development}

    \showmytikz{../shared/tikz/tikz-PM-tuckman-wheel}
\end{frame}

\begin{frame}{Tuckman's Phases of Team Development}
    
    \mytext{
  \begin{itemize}
      \item All stages repeat with each change in scope or team.
      \item Productive only in "performing" stage
      \item Accidents happen in "storming"
      \item Team building is a permanent effort
  \end{itemize}
  
  
  \vfill
   
  \parbox{0.5\linewidth}{
           Adding manpower to a project that is behind schedule delays it even more. 
           
           \flushright \textit{('Brook's Law')}
  }
  
  
  
    }
\end{frame}

\begin{frame}{Team Work and Diversity}
    
    \mytext{
    \begin{itemize}
        \item `There is no `I' in team
        \item `\textbf{t}oll, \textbf{e}in \textbf{a}nderer \textbf{m}acht's'
    \end{itemize}
    }
\end{frame}

\begin{frame}{Team Work and Diversity}
    \showmytikz{../shared/tikz/tikz-team-i}
\end{frame}

\begin{frame}{Team Work and Diversity}
    
    \mytext{
    \begin{itemize}
        \item What exactly is a team?
        \item Everybody equal, all different?
        \item Diversity in a team
        \item Coping with diversity is difficult
    \end{itemize}
    }
\end{frame}

\begin{frame}{Team Work and Diversity}
    
    \mytext{
    \begin{itemize}
        \item Don't surround yourself with yes-men 
        \item Teams should have complimentary qualities
        \item That is also a source of difficulty
        \item Different ways to communicate
        \item Different ways to cope with groups
        \item Different ways to cope with hierarchy
        \item ...
        \item Diversity is a plus!
    \end{itemize}    
    }
\end{frame}

\begin{frame}{Work Culture}
    
    \mytext{
    \begin{itemize}
        \item Meet online / chat rooms?
        \item Meet online / audio?
        \item Meet online / video?
        \item Meet IRL (in real life) - short
        \item Work together IRL
        \item Work together online
    \end{itemize}
    
    }
\end{frame}

\begin{frame}{PDCA (again)}

\showmytikz{../shared/tikz/tikz-PM-pdca-wheel-2}


    
\end{frame}

% \begin{frame}{Updated Module Plan}
% %\tikzset{every picture/.style={yscale=0.5}}
%     \begin{ganttchart}[y unit chart=0.6cm,
%     Milenew/.style={milestone/.append style={fill=red}}]{10}{17}
%       %\gantttitle{PM Module}{7} \\
%       \gantttitlelist{10,...,14}{1}\\
%       \gantttitlelist{24,...,31}{1} \\
%       %\ganttgroup{Participation}{7}{14} \\
%       %\ganttbar{Research existing project}{1}{3} \\
%       %\ganttmilestone{Presentation}{3} \\
%       %\ganttgroup{Own Project}{7}{14} \\
%       \ganttbar{Lectures}{10}{10} \\
%       \ganttbar[bar/.append style={fill=blue}]{Autonomous Work}{10}{14} \\
%       %\ganttbar{Identify Project}{7}{8} \\
%       %\ganttmilestone[Milenew]{One-pager Project Outline}{8} \\ %\ganttnewline 
%       \ganttbar[bar/.append style={fill=red}]{Prepare Presentations}{10}{14} \ganttnewline
%       \ganttmilestone{Presentations}{14} \ganttnewline
%       %\ganttmilestone{}{14} \\ %\ganttnewline 
%       \ganttbar{Write Documents}{10}{17} \ganttnewline
%       \ganttmilestone{Hand in Documents}{17}
%     \ganttvrule{\footnotesize{\textbf{July 14th}}}{14}
%     \ganttvrule{\footnotesize{Aug 1st}}{17}
%     \end{ganttchart}
% \end{frame}

\begin{frame}{Recommendations}
    
    \mytext{
    \begin{itemize}
        \item Use regular meetups / daily huddle
        \item Clarify roles / split requirements
        \item Assign tasks to individuals / \textbf{pairs}
        \item Agree on time table /dates
        \item Agree on technology / platform / storage
        \item Start early, prototype, etc.
        \item Schedule deadline before actual deadline
    \end{itemize}
    }
\end{frame}




\begingroup
    \ifdefined\setbeamercolor\setbeamercolor{background canvas}{bg=violet!20}\fi
    \begin{frame}{TASK}
    
    \mytext{
    \begin{itemize}
        \item For your proposal, create regular meetings
        \item Plan the tasks of your final assignement
    \end{itemize}
 
    
    }
\end{frame}
\endgroup